% ============================================================
% Conclusion Générale
% ============================================================
\chapter*{Conclusion Générale}
\addcontentsline{toc}{chapter}{Conclusion Générale}

Ce projet de fin d'études a été mené au sein de la société Amsys Consulting, dans le prolongement de son activité en intelligence artificielle et en ingénierie logicielle. La plateforme développée s'adresse aux opérateurs de télécommunications qui cherchent à prendre en charge automatiquement les demandes courantes de leurs abonnés sans renoncer au contrôle de ce que le système fait en leur nom. Elle se situe à l'intersection de trois domaines que le projet a dû articuler plutôt que juxtaposer : les systèmes conversationnels temps réel, la recherche d'information ancrée dans un corpus documentaire, et l'architecture logicielle appliquée à des opérations métier sensibles. Le cadre de sécurité et les pratiques d'observabilité de la société ont directement orienté plusieurs des choix présentés dans ce rapport.

Le travail réalisé a consisté à concevoir puis à construire cette plateforme dans son ensemble. Nous avons établi une architecture en ports et adaptateurs plaçant les règles métier hors d'atteinte des dépendances techniques, et découpé le domaine en contextes délimités communiquant par contrats explicites. Nous avons réalisé une chaîne de traitement vocale temps réel complète, dont chaque étape dispose d'une solution de repli éprouvée par injection de défaillance, et dont le budget de latence a été analysé étape par étape. Nous avons construit une couche agentique organisée autour de cinq agents spécialisés, dont le contexte est assemblé selon les besoins plutôt qu'accumulé, et dont les instructions sont dérivées des capacités réellement détenues, ce qui interdit qu'un agent soit invité à faire ce qu'il ne peut pas faire. Nous avons mis en place une chaîne documentaire dont la sélection procède par filtres successifs et dont la propriété première est de savoir se taire lorsque le corpus ne répond pas, propriété obtenue à la suite de mesures qui ont invalidé une approche initiale par seuil unique. Nous avons enfin réalisé la chaîne de décision, de politique et d'exécution qui protège les opérations sensibles, avec verdict motivé, exécution garantie unique et registre d'intégrité vérifiable, ainsi que les deux interfaces web et le dispositif d'observabilité qui rendent l'ensemble exploitable. Le système compte aujourd'hui vingt deux services conteneurisés, treize interfaces de domaine, soixante treize points d'accès applicatifs et trente cinq écrans, livrés au fil de cent douze incréments successifs.

Plusieurs perspectives se dégagent de ce travail, toutes réalisables à partir de l'état actuel du système. La première est la conduite d'une campagne de mesure instrumentée complète, avec exécutions répétées et analyse statistique, qui viendrait remplacer par des relevés expérimentaux les valeurs de latence reconstituées présentées dans ce rapport ; l'instrumentation nécessaire est déjà en place. La deuxième porte sur la recherche documentaire : l'asymétrie linguistique mise en évidence par les mesures justifierait des seuils différenciés par langue, ainsi qu'un élargissement du corpus permettant de vérifier que les réglages retenus sur un ensemble restreint se maintiennent à plus grande échelle. La troisième concerne l'exécution locale du modèle de raisonnement, écartée dans les conditions matérielles du projet mais redevenue envisageable sur une machine dotée d'un accélérateur graphique adapté, l'architecture ayant précisément été conçue pour que ce changement ne se propage pas au delà de la couche qui l'intègre. La quatrième consisterait à étendre la couverture métier en ajoutant de nouveaux domaines de spécialité, opération que la déclaration unique des domaines rend peu coûteuse. La cinquième, enfin, viserait l'ouverture de la plateforme à d'autres canaux que la voix et le texte, la couche de domaine étant déjà indépendante du canal d'entrée. Au delà du secteur des télécommunications, l'organisation retenue, qui associe une couche conversationnelle à une chaîne de décision déterministe et vérifiable, s'applique à tout domaine où un système automatisé doit agir sur des données engageant un client et rendre compte de ses décisions.
